\section{Introdução}

A atividade econômica está associada a externalidades ambientais, entre elas as emissões de dióxido de carbono (CO\textsubscript{2}), que contribuem para as mudanças climáticas. Os efeitos das mudanças climáticas já incluem danos à saúde, aos ecossistemas e à segurança alimentar e hídrica, com agravamento esperado no curto prazo. A redução das emissões integra os esforços globais de mitigação das mudanças climáticas \parencite{ipcc2023spm}.

A alocação dos esforços de mitigação depende de quais agentes são responsabilizados por quais emissões. Por isso, é relevante discutir os critérios utilizados para contabilizar as emissões e atribuir responsabilidade aos agentes econômicos. No âmbito do Protocolo de Kyoto, a contabilização das emissões nacionais segue uma base territorial, registrando-as no país em que ocorrem \parencite{unfccc2008kyotoManual}. A localização das fontes emissoras, contudo, não é o único critério possível de atribuição.

As mesmas emissões podem produzir distribuições setoriais distintas conforme o critério adotado. Pela produção, as emissões são contabilizadas na atividade em que ocorrem; pela demanda, são associadas ao setor do produto final, considerando as emissões diretas e indiretas necessárias ao atendimento de sua demanda final \parencite{marques2012,millerBlair2009}. As emissões da geração de eletricidade, por exemplo, podem ser registradas no setor elétrico ou atribuídas aos produtos finais cuja produção requer essa eletricidade. A mudança de critério preserva o total contabilizado, mas altera sua distribuição entre as atividades.

A escolha de quem responsabilizar pelas emissões envolve uma dimensão normativa. Identificar a atividade em que a emissão ocorre fornece uma informação relevante, mas não determina, por si só, como distribuir os esforços de mitigação. Em termos gerais, questões de justiça exigem reflexão e deliberação pública sobre os valores e os fins que devem orientar a vida em sociedade \parencite{sandel2015justica}. No caso das emissões, a comparação entre critérios de atribuição pode informar essa discussão, sem resolver a escolha normativa.

A literatura sobre emissões brasileiras oferece elementos para comparar atribuições entre setores, regiões e cadeias produtivas \parencite{imori2015,sanguinetAzzoni2024,montoyaEtAl2026}. A comparação dos totais atribuídos a cada atividade, entretanto, não mostra, isoladamente, quais atividades emissoras atendem à demanda final de quais produtos. Um volume elevado de emissões atribuídas pela demanda pode provir de poucas origens emissoras ou distribuir-se entre muitas. Da mesma forma, um grande emissor pode atender à demanda final de diferentes atividades.

Além dos totais, as diferenças entre as contas pela produção e pela demanda precisam ser examinadas em sua composição. Uma atividade pode apresentar uma diferença pequena entre as duas contas e, ao mesmo tempo, ter grandes volumes de emissões próprias atribuídas aos produtos finais de outras atividades e de emissões de outras atividades atribuídas aos seus produtos finais. Como esses volumes se compensam no saldo, a diferença entre as contas não revela sua magnitude nem sua distribuição. Identificar as atividades envolvidas em cada atribuição permite distinguir volumes concentrados em poucas relações de volumes distribuídos entre muitas.

Este trabalho pergunta: \textbf{quais relações intersetoriais compõem as diferenças entre as emissões atribuídas pela produção e pela demanda na economia brasileira de 2015, e como os volumes atribuídos se distribuem entre as atividades?} O objetivo é identificar as origens emissoras e os setores de produtos finais envolvidos nessas relações, caracterizar o volume e a concentração das atribuições e examinar as conexões entre as atividades. A análise busca, assim, complementar a comparação dos totais setoriais e discutir implicações para a coordenação da mitigação.

Para relacionar produção e demanda, utiliza-se a extensão ambiental do modelo insumo-produto \parencite{millerBlair2009}. A base econômica é a Matriz de Insumo-Produto (MIP) brasileira de 2015, com 67 atividades \parencite{ibge2015mip}, combinada com coeficientes setoriais de emissão de \textcite{sanguinetAzzoni2024}. A matriz de atribuição resultante é representada como uma rede dirigida e ponderada, adaptando a abordagem de \textcite{huang2019}. Cada nó representa uma atividade, e as ligações dirigidas associam origens emissoras a setores de produtos finais, com pesos correspondentes aos volumes de emissões atribuídos.

A análise distingue a distribuição dos volumes de emissões da presença de ligações entre atividades. Os volumes e a concentração são examinados na rede completa, antes da seleção de ligações pelo filtro de disparidade de \textcite{serrano2009}. O filtro identifica ligações destacadas na distribuição local dos pesos e produz uma rede reduzida, denominada \textit{backbone}, cuja conectividade é examinada. O número de ligações selecionadas não representa, por si só, a concentração dos volumes originais, razão pela qual as duas medidas são interpretadas separadamente.

O recorte abrange emissões da produção nacional associadas à demanda doméstica e às exportações, excluindo emissões estrangeiras. Os coeficientes ambientais compõem um cenário exploratório, com validação ambiental e monetária pendente, conforme detalhado na seção~\ref{sec:dados}.

Os resultados do cenário exploratório mostram uma rede quase completa: o número de ligações pouco diferencia as atividades, enquanto os pesos revelam diferenças de volume e concentração. A filtragem reduz o número de ligações, mas preserva proporções desiguais dos volumes de cada atividade. Esses resultados descrevem a estrutura do cenário analisado e não constituem um inventário validado das emissões brasileiras de 2015.

A discussão de política se concentra em conjuntos de setores e demandas que mereçam investigação conjunta. Intervenções nos processos produtivos e no uso de energia e materiais ao longo das cadeias constituem frentes complementares de mitigação \parencite{rissman2020}. A rede pode informar quais relações examinar, mas não estima a eficácia nem o custo de instrumentos de política.
